% Chapitre : méthodes de projection et banc de comparaison (axe 2).
% Sources : v2/docs/lab/reports/2026-11-1[45]-*, v2/methodes/out/*.json|log|csv, v2/papers/axe2-methodes/main.tex.
% Provenance : \mesure mesuré, \derive dérivé, \suppose supposé (macros de ../template/preamble.tex).
\begingroup\sloppy % chemins longs en \texttt ; groupe local au chapitre
\section{Méthodes de projection et banc de comparaison}\label{sec:methodes-ref}
\label{met:sec}

Ce chapitre définit chaque estimateur essayé pour prévoir l'effectif \enquote{libéral} (secteur privé,
lucratif ou non) des sages-femmes (SF), chirurgiens-dentistes (CD), pharmaciens (PH),
masseurs-kinésithérapeutes (MK) et infirmiers (IDE) par département, décrit le banc qui les compare, donne
les résultats et ce qui les limite. Le projet construit des \emph{estimateurs}, pas un nombre : chacun
passe par l'interface commune \texttt{fit(t, fit\_years, origin)} puis \texttt{predict(target)}
(\texttt{v2/sante/estimators.py}, classe \texttt{Estimator} ; adaptateur vers le jumeau dans
\texttt{v2/methodes/estimators.py}, classe \texttt{MethodEstimator}), et le banc les départage. Aucune
valeur de ce chapitre n'est \enquote{la} projection.

\paragraph{Abréviations des sources.} Les rapports cités sont dans \texttt{v2/docs/lab/reports/} :
\texttt{2026-11-14-axe2-methodes.md} (banc des variantes stock-flux), \texttt{2026-11-14-axe2-apc-sorties.md}
(APC), \texttt{2026-11-14-axe2-cohortes-bayesiennes.md} (C1, estimateur de cohortes ; à ne pas confondre avec
le critère d'acceptation C1 du jumeau ; C2), \texttt{2026-11-14-axe2-communes.md}
(K, K2), \texttt{2026-11-15-budget-erreur-c1-a2.md} (budget d'erreur), \texttt{2026-11-15-c1-intervalles.md}
(intervalles) et \texttt{2026-11-15-c1-gbm-global.md} (GBM, MLP). Les sorties machine sont dans
\texttt{v2/methodes/out/}. Les décisions sont dans \texttt{v2/docs/lab/decisions.md}, les questions
ouvertes dans \texttt{v2/docs/lab/open.md}.

\subsection{Notation et identité comptable}
\label{met:notation}
Une strate $c$ est un triplet (profession, département $d$, sexe) ; dans ce rapport, \emph{cellule} désigne
un couple profession $\times$ origine du banc, et \emph{couple} un département $\times$ cellule ; $a$ indexe la tranche d'âge
quinquennale ($<25$, 25--29, \dots, 65+) ; $t$ l'année. Le RPPS (DREES, publication du 2026-07-02) ne
donne que des stocks $N_{c,a}(t)$, jamais les entrées ni les sorties
(\texttt{v2/papers/axe2-methodes/main.tex}, section Données). Toutes les méthodes stock-flux partagent
l'identité
\begin{equation}
  N_c(t+1) = s_c \odot A\big(N_c(t)\big) + I_c,
  \label{met:eq:stockflux}
\end{equation}
où $A$ est l'opérateur de vieillissement, $s$ la rétention nette et $I$ le flux net d'entrées. Avec des
âges uniformes dans la tranche, $A$ fait passer $1/5$ de chaque tranche à la suivante :
$A(N)_a = 0{,}8\,N_a + 0{,}2\,N_{a-1}$ pour les tranches intérieures, $A(N)_{<25} = 0{,}8\,N_{<25}$ et
$A(N)_{65+} = N_{65+} + 0{,}2\,N_{60\text{--}64}$ aux bornes (\texttt{v2/sante/model.py}, fonction
\texttt{\_inflow} ; vérifié à $1{,}6\times10^{-14}$\mesure{} près,
\texttt{v2/papers/axe2-methodes/main.tex}, annexe). Le ratio de cohorte observé est
$r_{c,a}(\tau) = N_{c,a}(\tau)/A(N_c(\tau-1))_a$ : il mêle cessations, reprises, changements de secteur et
déménagements. La sortie nette reconstruite $A(N_c(\tau-1)) - N_c(\tau)$ est négative dans 27 à 47\,\%\mesure{} des
strates $\times$ tranches de 35 ans et plus (\texttt{2026-11-14-axe2-apc-sorties.md}) ; c'est pourquoi aucune méthode retenue ne
modélise des cessations brutes.

\subsection{Famille stock-flux déterministe}
\label{met:stockflux}
Les sept variantes suivantes gardent l'équation~\eqref{met:eq:stockflux} et ne changent que $A$, le partage
d'information entre territoires ou la dépendance de $s$ au temps. Elles ont été inscrites au registre
(\texttt{v2/validation/registry.py}) avant leur premier passage ; dérivation déposée sous l'empreinte
SHA-256 \texttt{1c324f5795b0} (\texttt{v2/papers/axe2-methodes/main.tex}, § Prédictions).

\paragraph{A0, stock-flux à taux locaux (référence).} Estimateur \texttt{jumeau\_taux\_locaux}, jumeau
\texttt{twin\_of(StockFlowLocal)}. $s_{c,a} = \min(1, r_{c,a})$, où $r$ est un ratio de sommes sur la
fenêtre, rétréci vers le national (\texttt{v2/sante/model.py}, fonction \texttt{fit\_rates}) :
\begin{equation}
  r_{c,a} = \frac{\text{num}_{c,a} + k\, r_{\text{nat},a}}{\text{den}_{c,a} + k}, \qquad
  \text{num}_{c,a} = \textstyle\sum_\tau N_{c,a}(\tau), \quad
  \text{den}_{c,a} = \textstyle\sum_\tau A\big(N_c(\tau-1)\big)_a, \qquad k = 200\suppose,
  \label{met:eq:shrink}
\end{equation}
avec $r_{\text{nat},a}$ le même ratio sur la France (par profession $\times$ sexe) ; $k$ s'exprime en
personnes exposées (poids de l'a priori national dans le dénominateur),
et $I_{c,a}$ le flux net moyen de la fenêtre dans les tranches où le ratio brut dépasse~1. C'est la seule
variante confirmatoire du protocole N1 (\texttt{decisions.md}, entrée \enquote{Axe 3, N1 : grille,
métrique, seuil et variante par défaut}). Elle est servie par le jumeau si \texttt{JUMEAU\_MODELE=A0}.

\paragraph{A1a, âges interpolés.} Estimateur \texttt{axe2\_ages\_interp}. Chaque année, chaque tranche est
désagrégée en âges simples par une densité linéaire par morceaux entre milieux de tranches, tronquée à 0 et
remise à l'échelle pour conserver la masse de la tranche ; on vieillit d'un an puis on réagrège. Le même
opérateur sert à l'ajustement et à la projection. Hypothèses : tranche $<25$ = 20--24 ans\suppose{},
tranche 65+ de largeur 10 ans\suppose.

\paragraph{A1b, âges simples.} Estimateur \texttt{axe2\_ages\_simples}. L'état est
$N[c, x]$, $x = 20, \dots, 74+$, désagrégé une seule fois à l'origine ; les cohortes vieillissent
exactement d'un an. $s$ et $I$ restent calibrés par tranche ; les entrées sont réparties uniformément dans
la tranche\suppose.

\paragraph{A2, shrinkage hiérarchique.} Estimateur \texttt{axe2\_shrink\_hier}. L'équation
\eqref{met:eq:shrink} est appliquée à deux niveaux :
\begin{equation}
  r_{\text{dép}} = \frac{\text{num} + k\, r_{\text{rég}}}{\text{den} + k}, \qquad
  r_{\text{rég}} = \frac{\text{num}_{\text{rég}} + k\, r_{\text{nat}}}{\text{den}_{\text{rég}} + k},
  \qquad k = 200\suppose,
\end{equation}
poids repris de la v1 et de l'observatoire interne, non réglé sur le banc.

\paragraph{A3, rétention des 60 ans et plus en tendance.} Estimateur \texttt{axe2\_retention60}. Pour
$a \ge 60$ :
$s_{c,a}(t) = \min\!\big(1,\ s_{c,a} + \lambda\, b_{p,\text{sexe},a}\,(t - \bar t)\big)$, où $b$ est la
pente nationale des ratios annuels (moindres carrés ordinaires) et $\bar t$ le milieu de la fenêtre ;
$\lambda = 1$ dans la variante inscrite. Le levier \texttt{methodes\_sorties.tendance\_60plus} porte
$\lambda$ (0 par défaut).

\paragraph{A123.} Estimateur \texttt{axe2\_combine} : A1a + A2 + A3, combinaison pré-spécifiée.

\paragraph{B4, schéma de l'observatoire (diagnostic, non candidat).} Estimateur \texttt{axe2\_obs\_r12}.
Moins de 35 ans : $r = 1$ et entrées = résidus moyens sur 5 ans ; ailleurs $r$ rétréci
département $\to$ région $\to$ national et borné à $[0{,}5\,;\,1{,}2]$, sans entrée absolue. Déclaré
diagnostic avant le passage.

\paragraph{Variantes jugées sans passage au banc.} C5 (la part régionale des étudiants prédit-elle la part
régionale des entrées nettes ?) est un diagnostic. A4 (logit conditionnel du lieu d'installation, avec effets
fixes départementaux) reproduit exactement les parts observées\derive{} et n'est utile qu'avec un attribut
daté : en attente. A5 (bayésien hiérarchique, espace-état) donne au département $N \ge 30$ une moyenne a
posteriori proche de A2\derive{} ; il a été réalisé de fait par l'APC et C1. A6 (microsimulation) a la
même espérance que A0 par linéarité de $N(t+1) = S\,A\,N(t) + I$\derive{} : tuée pour N1. Source :
\texttt{2026-11-14-axe2-methodes.md} et la dérivation \texttt{v2/docs/lab/derivations/2026-11-14-axe2-methodes.md}.

\subsection{Famille bayésienne sur la structure par cohortes}
\label{met:bayes}

\paragraph{APC, âge-période-cohorte gaussien de la rétention nette.} Estimateur \texttt{axe2\_apc},
module \texttt{apc\_sorties}, code \texttt{v2/methodes/apc.py}. Par profession $\times$ sexe, pour chaque
observation $i = (d, a, \tau)$ dans les tranches où la rétention nationale de la fenêtre est $\le 1$ :
\begin{align}
  y_i &= \log\frac{N_i + 1/2}{B_i + 1/2} \sim \mathcal N\!\left(\eta_i,\ \sigma_e^2/B_i + \omega^2\right),\\
  \eta_i &= \mu_a + \pi_{g(a),\tau} + \gamma_{\tau-a} + v_{\text{rég}(d),a} + w_{d,a},
\end{align}
où $B_i$ est l'effectif exposé (vieillissement uniforme), $\mu_a \sim \mathcal N(0,10^2)$, $\pi_g$ une marche
aléatoire d'ordre 1 par groupe d'âge ($<35$, 35--59, $\ge 60$), $\gamma$ une marche aléatoire d'ordre 1 sur
la cohorte indexée en années~\cite{b-schmid2007bamp}, $v \sim \mathcal N(0, \tau_r^2)$ et
$w \sim \mathcal N(0,\tau_d^2)$. Les hyperparamètres
$\theta = (\sigma_\pi, \sigma_\gamma, \tau_r, \tau_d, \sigma_e, \omega)$ maximisent la vraisemblance
marginale (Bayes empirique de type II, L-BFGS-B) ; sachant $\hat\theta$, la loi a posteriori des 500 à
700\mesure{} effets est exactement gaussienne (factorisation de Cholesky dense). En projection, les marches
sont prolongées à plat et $s = \min(1, e^{\eta})$ ; les entrées restent celles de A0. Levier :
\texttt{apc\_sorties.periode\_60plus}. Source : \texttt{2026-11-14-axe2-apc-sorties.md}.

\paragraph{C1, cohortes par composantes, $\theta$ fixé (servi par le jumeau).} Estimateur
\texttt{axe2\_cohortes\_c1}, classe \texttt{CohortesEstimator}, code \texttt{v2/methodes/cohortes.py}.
Sorties : l'APC ci-dessus. Entrées : dans les tranches où la rétention nationale dépasse 1, le flux net
$e_{d,\tau} = N(a,\tau) - s_d\,B(a,\tau-1)$ suit
\begin{align}
  M_\tau &= \textstyle\sum_d e_{d,\tau} = L_\tau + u_\tau, \qquad L_\tau = L_{\tau-1} + v_\tau, \\
  e_{d,\tau} &= \pi_d\, M_\tau\,\big(1 + \delta_{\text{rég}(d)} + \zeta_d\big) + \varepsilon_{d,\tau},
\end{align}
avec $u \sim \mathcal N(0,\sigma_u^2\kappa)$, $v \sim \mathcal N(0,\sigma_v^2\kappa)$, $\kappa$ la moyenne de
$|M_\tau|$, $\pi_d$ la part du département dans le stock de la tranche à l'origine (fixée),
$\delta \sim \mathcal N(0,\tau_r^2)$, $\zeta \sim \mathcal N(0,\tau_d^2)$,
$\varepsilon \sim \mathcal N(0, \sigma_\varepsilon^2\pi_d\kappa + \omega_e^2)$. Le niveau local $L$ est
filtré par Kalman ($\sigma_u, \sigma_v$ au maximum de vraisemblance) ; la prévision est
$M_{T+h} \sim \mathcal N\big(L_{T|T},\ P_{T|T} + h\,\sigma_v^2 + \sigma_u^2\big)$, tirée 200 fois. La forme
des variances (proportionnelles à $\kappa$) a été corrigée après un pilote qui donnait des variances
relatives explosives ($\sigma_u \approx 1{,}3$\mesure, $\sigma_v \approx 1{,}2$\mesure{} pour les SF) ; ce
pilote avait vu deux cellules cibles (\S\ref{met:limites}). Source :
\texttt{2026-11-14-axe2-cohortes-bayesiennes.md}.

C1 est \textbf{le modèle servi par le jumeau} (carte et assistant) depuis l'entrée \enquote{C1 servi par la
carte et l'assistant à la place de A0} de \texttt{decisions.md} : sa moyenne a posteriori remplace la
rétention et les entrées de base de A0 pour les cellules privées, au point où le jumeau calcule ses taux,
par un cache de l'historique (\texttt{v2/methodes/c1\_jumeau.py}, fonction \texttt{install}) ; les taux
sont précalculés dans \texttt{v2/methodes/out/c1\_taux/}. Le salariat public garde les taux de A0. La
prévision C1 origine 2026 est déposée pour le score 2027
(\texttt{v2/validation/depot/axe2\_cohortes\_c1\_2026-11-14.parquet}, SHA-256 \texttt{689b230b5b50},
\texttt{decisions.md}).

\paragraph{C2, C1 avec $\theta$ propagé.} Classe \texttt{CohortesLaplace} : approximation de Laplace sur
$\log\theta$, $\log\theta \sim \mathcal N(\log\hat\theta, H^{-1})$ avec $H$ la hessienne numérique, puis
effets gaussiens exacts sachant chaque $\theta_k$ tiré. Coût double\mesure{} de C1.

\paragraph{Écartés a priori dans cette famille.} Un Poisson sur sorties reconstruites $D = A(N) - N$ exige
$D \ge 0$, faux dans 27 à 47\,\%\mesure{} des strates $\times$ tranches $\ge 35$ ans
(\texttt{v2/docs/lab/derivations/2026-11-14-axe2-apc-sorties.md}). Un APC d'ordre 2, qui prolongerait la
pente, est écarté parce que A3 a échoué. Le MCMC complet (PyMC, NumPyro) a été tué pour l'étape 2 puis
rouvert après l'échec de la calibration par simulation de $L_T$ ; il n'est pas fait (\texttt{open.md}).

\subsection{Échelle communale : K à K5}
\label{met:communes}
Toutes sont jugées sur CartoSanté (libéraux au 31/12 ; IDE, MK, CD ; origines 2014--2020 ; $h = 5$) par le
protocole N4 et ses suites N10--N12 ; critère : bande à 80\,\% couvrant 75 à 85\,\% des cas dans chaque
classe de taille (\texttt{2026-11-14-axe2-communes.md}). C1 tel quel à la commune est tué a priori (pas de
sexe, tranches CartoSanté inégales, densité de paramètres).
\begin{description}
  \item[K, cohortes communales hiérarchiques.] Vieillissement
    $B(a,\tau) = (1 - f_a)\,N(a,\tau-1) + f_{a-1}\,N(a-1,\tau-1)$,
    $f = (1/15, 1/10, 1/5, 1/5, 0)$ (tranche $<40$ = 25--39 ans\suppose) ; rétention rétrécie
    national $\to$ région $\to$ département $\to$ EPCI $\to$ commune ($k = 200$\suppose) ; entrées = niveau
    local départemental (Kalman) $\times$ part communale $q_c$ rétrécie ($k_q = 20$\suppose) ; variance
    $\alpha_h\mu + \beta_h\mu^2$ mesurée en pseudo-prévisions dans la fenêtre ; binomiale négative si
    $\mathrm{Var} > \mu$, Poisson sinon ; bande P10--P90.
  \item[K2.] Point = part communale tenue $\times$ total départemental de K ; bande = quantiles 10 et 90\,\%
    de $\log\frac{\text{réel}+1/2}{\text{prédit}+1/2}$ par classe de taille.
  \item[K3.] Point de K2 ; $X_c = (p_c + 1/2)\exp(\varepsilon_d + \eta_c) - 1/2$, tronqué à 0, où
    $\varepsilon_d$ est rééchantillonné parmi les erreurs log des totaux départementaux de K (commun aux
    communes d'un département) et $\eta_c$ est une erreur locale par classe ; pseudo-origines
    $o - 5 - \ell$, $\ell = 0,\dots,4$ ; $S = 200$ ; bassins de vie et EPCI par somme des tirages.
  \item[K4.] K3 avec récence ($\ell \in \{0, 1\}$) et tirage conjoint de $\varepsilon_d$, $\eta_c$ dans la
    même pseudo-origine.
  \item[K5.] $\log\frac{n_c + 1/2}{p_c + 1/2} = b_p + \sigma_{D,p}\,z_d + \sqrt{\phi_p/(n^0_c + 1/2) + \kappa_p}\;z_c$ ;
    $\sigma_D$ par écart-type robuste ($1{,}4826\times$ MAD), $b$ = médiane, $\phi, \kappa$ par maximum de
    vraisemblance gaussien (erreurs écrêtées aux centiles 1 et 99) ; $\log\sigma_D, \log\phi, \log\kappa$
    rétrécis de moitié vers la moyenne des trois professions (poids 0,5\suppose).
\end{description}

\subsection{Familles hors stock-flux jugées a priori}
\label{met:horssf}
Source : dérivation \texttt{v2/docs/lab/derivations/2026-11-14-axe2-familles-sd-abm.md} (abrégée D-sd
dans \texttt{decisions.md}) et \texttt{2026-11-15-budget-erreur-c1-a2.md}.
\begin{description}
  \item[D1, boucle de dynamique des systèmes (IDE, 45--54 ans).]
    $\log s_{d,45\text{--}54}(t) = \log s_{d,45\text{--}54} + \beta\,\log\big(\rho_d(t)/\rho_d(t_0)\big)$, avec
    $\rho_d$ la densité de l'état simulé et $\beta \sim \mathcal N(0{,}11\,;\,0{,}03^2)$ (protocole N9). En
    attente, comme scénario de stress seulement.
  \item[D2, agents à règles calées.] Règles individuelles calées sur les stocks : non identifiables sans panel
    individuel ; tuée.
  \item[D3, agents à règles estimées.] Choix de commune selon les probabilités d'un logit estimé ; la
    régulation devient un ensemble de choix qui se ferme. En attente du logit.
  \item[D4, multi-états de Markov privé $\leftrightarrow$ public $\leftrightarrow$ sortie.] Matrice de
    transition annuelle $P$ par tranche $\times$ sexe $\times$ département ; seules les transitions nettes
    sont identifiables. Promue, non construite (\texttt{open.md}).
  \item[Modèle gravitaire de mobilité.] Rejeté par le budget d'erreur (\S\ref{met:budget}) : persistance
    de l'écart local $r$ = 0,22--0,24, non établie au seuil de 0,3, et aucun flux origine--destination ouvert.
\end{description}

\subsection{Apprentissage automatique global : GBM et MLP}
\label{met:gbm}
Code \texttt{v2/methodes/gbm.py} ; dérivation inscrite avant le banc
(\texttt{v2/docs/lab/derivations/2026-11-15-c1-gbm-global.md}, SHA-256 \texttt{8e35043aecb2}). Un seul
modèle pour toutes les professions, départements et tranches, à un pas et itéré jusqu'à la cible. Unité :
strate $\times$ tranche (profession $\times$ département $\times$ sexe $\times$ tranche). Base sans apprentissage
$b = 0{,}8\,n(a) + 0{,}2\,n(a-1)$ (65+ : $n + 0{,}2\,n(a-1)$). Cible
$y = \big(n(a,t+1) - b(t)\big)/N(d,t)$, la contribution de la cellule à la croissance relative du
département. Vingt variables : retards de $y$ et du ratio brut, croissances départementales à 1, 2 et 3 ans,
croissance nationale, parts des 60+ et des $<35$, population (log), part des 60+ dans la population,
densité, et profession, sexe, tranche, région en catégories (\texttt{FEATURES}, \texttt{gbm.py}).
\begin{description}
  \item[GBM.] \texttt{HistGradientBoostingRegressor}, 300 itérations, taux 0,05, 15 feuilles, au moins 50
    observations par feuille, pénalité L2 = 1, sans arrêt précoce ; réglages fixés a priori\suppose.
  \item[MLP.] \texttt{MLPRegressor} à couches (32, 16), ReLU, $\alpha = 10^{-3}$, arrêt précoce sur 10\,\% des
    lignes d'entraînement, environ 1\,500\derive{} poids.
\end{description}
Le panel est coupé à l'origine ; un contrôle de fuite (perturbation des années postérieures à 2019) donne
un écart maximal de prédiction de 0,0\mesure{} pour les deux modèles (\texttt{v2/methodes/out/c1\_gbm.log}).
Ce contrôle n'est joué qu'à l'origine 2019 (dentistes ; \texttt{v2/methodes/c1\_gbm.py}, fonction
\texttt{leak\_check}) : il ne vaut que pour les origines $\le 2019$. Pour les origines 2020 et 2021, seules
les assertions du code (\texttt{gbm.py} : aucune donnée postérieure à l'origine) s'appliquent.
Un modèle direct à $h = 5$ est infaisable sans fuite aux origines 2016--2017\derive.

\subsection{Intervalles de prédiction de C1 : conforme et MinT}
\label{met:intervalles}
Code \texttt{v2/methodes/c1\_intervalles.py}, hors protocole N6 (gelé, non modifié). On note $m$ et $s$ la
médiane et l'écart-type des log-tirages de C1 et $z = (\log y - m)/s$ le score normalisé d'une cible
passée.
\begin{description}
  \item[C1-conf-h5.] Prédiction conforme au même horizon : scores des couples $(o', h' = 5)$ avec
    $o' + 5 \le o$ ; prédictive $m + s\,z$. Seule l'origine 2021 est calculable ($o' = 2016$).
  \item[C1-conf.] Même chose sur tous les horizons $h' \le 5$ avec $o' + h' \le o$ ; origines 2017--2021.
  \item[C1-conf-hg (post hoc).] $z$ ramené à $h = 5$ par $\mathrm{sd}(z \mid h') \propto h'^{\gamma}$, ajusté
    sur le passé ; écrit après avoir vu que $\mathrm{sd}(z)$ croît avec $h'$.
  \item[C1-MinT.] Tirages réconciliés département $\to$ région $\to$ France par MinT-WLS ($W$ = variances des
    tirages de base), chaque tirage projeté ; la région et la France sont ajustées en relabellisant le
    département.
  \item[C1-MinT-conf, C1-MinT-conf-hg.] MinT puis conforme sur ses propres scores passés.
\end{description}
Les références de méthode (prédiction conforme, conforme adaptatif, MinT) sont données \enquote{de mémoire}
dans \texttt{2026-11-15-c1-intervalles.md} et n'ont pas été lues en texte intégral : la règle du projet
(\texttt{decisions.md}) les exclut de la bibliographie de ce rapport.

\subsection{Ce que la carte affiche}
\label{met:servi}
La carte sert la moyenne a posteriori de C1. La fourchette affichée n'est pas tirée de C1 : pour une valeur
servie $v$ à l'horizon $h$, c'est $v\,(1 \pm k_p\,h)$, linéaire jusqu'en 2050
(\texttt{v2/prototype/index.html} l.~225 ; \texttt{v2/prototype/server.py} l.~158--180). La pente $k_p$ est
celle de l'erreur relative absolue du total \emph{national} de C1 au banc du 2026-11-14, ajustée par moindres
carrés sans constante sur $h$, et inclut déjà le facteur 1,6 (\texttt{v2/methodes/c1\_jumeau.py}, fonction
\texttt{calcule\_pente}, l.~111) : \texttt{v2/methodes/out/c1\_pente.json} donne CD 0,0172\mesure,
IDE 0,0132\mesure, MK 0,0079\mesure, PH 0,0037\mesure, SF 0,0172\mesure. La même pente nationale est
appliquée telle quelle à chaque territoire. Cette fourchette n'est donc calibrée ni au département ni au-delà
de $h = 10$. La phrase de \texttt{decisions.md} et de \texttt{2026-11-15-c1-intervalles.md} selon laquelle
\enquote{la bande affichée reste celle des tirages de C1} ne décrit pas le code (annexe,
tableau~\ref{met:tab:inc}).

\subsection{Protocole du banc N1}
\label{met:protocole}
Le protocole est gelé sous l'empreinte \texttt{eb527ecb06a0} (\texttt{decisions.md}, entrée N1 ; code
\texttt{v2/validation/forecast.py}).
\begin{itemize}
  \item \textbf{Origines et fenêtres.} Origines glissantes ; fenêtre d'ajustement depuis 2012 (kinés : 2017,
    rupture ADELI $\to$ RPPS), d'au moins 5 ans (\texttt{MIN\_FIT = 5}). À $h = 5$ : CD, PH et SF 2016--2021
    (6 origines chacune), IDE 2016--2018 (3, série libérale arrêtée en 2023), MK 2021 (1), soit
    \textbf{22 cellules} profession $\times$ origine\mesure{} et \textbf{2\,056 couples} département
    $\times$ cellule\mesure{} (\texttt{v2/methodes/out/budget\_erreur.json}, clé
    \texttt{couples\_dep\_cellule}).
  \item \textbf{Horizons.} $h \in \{1, 3, 5, 8, 10\}$ (\texttt{HORIZONS}) ; $h = 5$ est le seul horizon
    commun aux cinq professions sur plusieurs origines.
  \item \textbf{Métrique principale unique.} MAPE par département à $h = 5$, départements d'au moins 30
    professionnels à l'origine (\texttt{MIN\_N = 30}), moyenne à poids égal sur les cellules :
    \begin{equation}
      \text{MAPE} = \frac{1}{|\mathcal C|}\sum_{(p,o)\in\mathcal C}\ \frac{1}{|D_{p,o}|}\sum_{d \in D_{p,o}}
      100\,\frac{|\hat N_{p,d}(o+5) - N_{p,d}(o+5)|}{N_{p,d}(o+5)} .
    \end{equation}
    WAPE, biais national signé et $h \ne 5$ sont secondaires.
  \item \textbf{Critère.} La borne haute de l'IC 95\,\% de $\text{MAPE(candidat)} - \text{MAPE(meilleure
    naïve)}$ doit être $< 0$. Naïves : persistance $\hat N(o+h) = N(o)$ et tendance linéaire ; la meilleure est
    choisie après coup, ce qui défavorise le candidat.
  \item \textbf{Bootstrap.} Rééchantillonnage des départements, 2\,000 tirages, graine 20261114
    (\texttt{SEED}, \texttt{BOOT}) ; mêmes tirages pour toutes les colonnes d'un banc. Le même tirage de
    départements sert pour toutes les cellules (\texttt{forecast.py} l.~58--65), ce qui conserve la
    corrélation entre origines d'un même département.
  \item \textbf{Périmètre confirmatoire.} Privé (variante centrale de l'axe 4) ; pour les IDE, libéraux du
    RPPS, l'âge des salariés privés manquant. Périmètre secondaire \enquote{ancien} : modes libéral et mixte.
  \item \textbf{Limite $h \to 0$.} Garde-fou unilatéral : à $h = 1$, la MAPE du candidat ne doit pas
    dépasser celle de la persistance de plus d'un point (\texttt{limite\_h1\_ok}, \texttt{forecast.py}).
\end{itemize}
Le protocole probabiliste N6 (SHA-256 \texttt{d1cc665cb864}) juge le CRPS sur $\log N$~\cite{b-gneiting2007jasa}
contre une naïve probabiliste (prédiction $\times\,e^{\varepsilon}$, $\varepsilon$ tiré parmi au moins 20
erreurs log passées dont la cible précède l'origine) ; seule l'origine 2021 (SF, CD, PH ; 279 couples\mesure)
est jugeable. Ce protocole juge A0 ; C1 y a été passé ensuite (\S\ref{met:resultats}). La calibration par
simulation suit~\cite{b-talts2018sbc}.

\subsection{Résultats du banc N1}
\label{met:resultats}

\paragraph{Table de référence : banc du 2026-11-15.} Le tableau~\ref{met:tab:ref15} est la référence de ce
rapport. Il réunit deux recalculs du 15 novembre sur la table privée actuelle : \texttt{methodes.budget\_erreur}
(22 réajustements de C1 et A2 ; \texttt{v2/methodes/out/budget\_erreur.json}) et \texttt{methodes.c1\_gbm}
(A2, naïves, GBM, MLP ; \texttt{v2/methodes/out/c1\_gbm\_verdicts.json}, \texttt{c1\_gbm\_cellules.csv}).
Les deux recalculs donnent la même MAPE pour A2 (8,60\,\%\mesure), et les MAPE par profession d'A2
concordent au centième.

\begin{table}[ht]
\centering\small
\resizebox{\textwidth}{!}{%
\begin{tabular}{@{}lrrrrrrl@{}}
\toprule
Estimateur & Total (IC 95\,\%) & CD & IDE & MK & PH & SF & Source \\
 & 22 cellules & 6 & 3 & 1 & 6 & 6 & \\
\midrule
\textbf{C1} (servi) & \textbf{7,76}\mesure{} $[7{,}12\,;\,8{,}40]$ & 8,03 & 4,93 & 6,51 & 4,91 & 11,96 & \texttt{budget\_erreur.json} \\
A2 & 8,60\mesure{} $[7{,}80\,;\,9{,}47]$ & 8,66 & 5,72 & 6,62 & 5,06 & 13,87 & idem et \texttt{c1\_gbm\_verdicts.json} \\
tendance (naïve) & 9,32\mesure & 8,70 & 7,62 & 9,12 & 5,85 & 14,29 & \texttt{c1\_gbm\_verdicts.json} \\
persistance (naïve) & 9,94\mesure & 8,22 & 10,84 & 15,48 & 5,67 & 14,54 & \texttt{c1\_gbm\_cellules.csv}\derive \\
GBM global & 10,15\mesure & 8,45 & 7,27 & 11,37 & 6,00 & 17,25 & \texttt{c1\_gbm\_verdicts.json} \\
MLP global & 20,52\mesure & 23,23 & 14,34 & 60,13 & 11,75 & 23,07 & idem \\
\bottomrule
\end{tabular}}
\caption{Banc N1, périmètre privé, MAPE par département à $h = 5$ (\%), table privée du 2026-11-15. Les
valeurs par profession sont mesurées\mesure{} sauf la persistance, moyennée par profession à partir des 22
lignes cellule de \texttt{c1\_gbm\_cellules.csv}\derive. Les IC sont des IC bootstrap \emph{du niveau} de
MAPE, à ne pas comparer aux IC d'écarts du tableau~\ref{met:tab:ref14}. A0, A1a, A1b, A3, A123, B4, APC
et C2 n'ont pas été recalculés sur cette table.}
\label{met:tab:ref15}
\end{table}

La MAPE totale est la moyenne à poids égal des 22 cellules ; elle se retrouve à partir des MAPE par
profession pondérées par le nombre de cellules :
$(6\times8{,}03 + 3\times4{,}93 + 6{,}51 + 6\times4{,}91 + 6\times11{,}96)/22 = 7{,}76$\derive{} pour C1,
10,15\derive{} pour le GBM et 8,61\derive{} pour A2 (8,60\mesure{} avant arrondi des MAPE par profession).
L'écart C1 $-$ A2 vaut $-0{,}84$\derive{} point. Le GBM est à 10,15 contre 8,60 pour A2 :
$+1{,}55$\mesure{} point, IC 95\,\% $[+0{,}62\,;\,+2{,}75]$\mesure{} (IC 97,5\,\% de Bonferroni
$[+0{,}53\,;\,+2{,}91]$\mesure) ; le MLP à $+11{,}92$\mesure{} $[+10{,}54\,;\,+13{,}86]$\mesure. Les deux sont
rejetés ; le GBM fait jeu égal avec A2 pour les CD ($-0{,}21$\mesure, $[-0{,}85\,;\,+0{,}43]$\mesure) et
les PH ($+0{,}94$\mesure, $[-0{,}48\,;\,+3{,}48]$\mesure) et perd ailleurs.

\paragraph{Le banc du 2026-11-14.} Le tableau~\ref{met:tab:ref14} donne le premier passage, sur la table
privée d'origine (\texttt{v2/methodes/out/verdicts\_prive.json} pour les variantes stock-flux,
\texttt{coh\_verdicts.json} pour APC, C1 et C2). Ses valeurs pour A2 et C1 (8,57 et 7,73) diffèrent de la
référence\footnote{Banc du 2026-11-14 : A2 8,57\,\%\mesure{} (\texttt{verdicts\_prive.json}), C1
7,73\,\%\mesure{} (7,7315, \texttt{coh\_verdicts.json}) ; banc du 2026-11-15 : 8,60\mesure{} et
7,76\mesure. Les deux viennent du même protocole N1 (mêmes 22 cellules, 2\,056 couples, même bootstrap).
Cause : la table privée des sages-femmes a été modifiée entre les deux bancs, jusqu'à 4,5\mesure{}
personnes par département ; les quatre autres professions sont identiques \enquote{à 0 près}
(\texttt{2026-11-15-budget-erreur-c1-a2.md}, § Ce qui pourrait rendre ce résultat faux). Contrôle : la
MAPE SF d'A2 passe de 13,73\mesure{} à 13,87\mesure, les autres professions d'A2 sont inchangées, et
$0{,}14 \times 6/22 = 0{,}038$\derive{} point, soit 8,57 $\to$ 8,60. Les naïves bougent dans le même
sens : tendance 9,30 $\to$ 9,32\mesure, persistance 9,92 $\to$ 9,94\mesure.}. Les écarts entre variantes y
restent valides à l'intérieur de ce banc, mais aucun de ses nombres ne doit être mélangé avec ceux du
tableau~\ref{met:tab:ref15}.

\begin{table}[ht]
\centering\small
\resizebox{\textwidth}{!}{%
\begin{tabular}{@{}lrrlrl@{}}
\toprule
Variante & MAPE & $\Delta$ vs tendance (IC 95\,\%) & Issue N1 & $\Delta$ vs A0 (IC 95\,\%) & Verdict \\
\midrule
A0 & 8,68 & $-0{,}63$ $[-0{,}97\,;\,-0{,}27]$ & gagne & --- & référence \\
A2 & 8,57 & $-0{,}74$ $[-1{,}07\,;\,-0{,}40]$ & gagne & $-0{,}11$ $[-0{,}21\,;\,-0{,}02]$ & bat A0 \\
A1b & 8,70 & $-0{,}60$ $[-0{,}96\,;\,-0{,}24]$ & gagne & $+0{,}03$ $[-0{,}07\,;\,+0{,}11]$ & égalité \\
A123 & 8,74 & $-0{,}56$ $[-0{,}92\,;\,-0{,}21]$ & gagne & $+0{,}06$ $[-0{,}07\,;\,+0{,}18]$ & égalité \\
A1a & 8,77 & $-0{,}53$ $[-0{,}89\,;\,-0{,}17]$ & gagne & $+0{,}09$ $[+0{,}01\,;\,+0{,}18]$ & échoue de peu \\
A3 & 9,00 & $-0{,}30$ $[-0{,}67\,;\,+0{,}08]$ & égalité & $+0{,}33$ $[+0{,}23\,;\,+0{,}42]$ & échoue \\
B4 (diagnostic) & 8,32 & $-0{,}98$ $[-1{,}37\,;\,-0{,}59]$ & gagne & $-0{,}35$ $[-0{,}65\,;\,-0{,}08]$ & non candidat \\
APC & 8,74 & $-0{,}56$ $[-0{,}91\,;\,-0{,}20]$ & gagne & $+0{,}07$ $[-0{,}03\,;\,+0{,}16]$ & égalité \\
C1 & 7,73 & $-1{,}57$ $[-2{,}17\,;\,-0{,}99]$ & gagne & $-0{,}95$ $[-1{,}51\,;\,-0{,}47]$ & bat A0 \\
C2 & 7,73 & $-1{,}57$ $[-2{,}17\,;\,-0{,}99]$ & gagne & $-0{,}95$ $[-1{,}51\,;\,-0{,}47]$ & bat A0, coût $\times2$ \\
tendance & 9,30 & --- & --- & --- & naïve \\
persistance & 9,92 & --- & --- & --- & naïve \\
\bottomrule
\end{tabular}}
\caption{Banc N1, périmètre privé, $h = 5$, table privée du 2026-11-14 ; toutes les valeurs sont
mesurées\mesure. Sources : \texttt{verdicts\_prive.json} (A0 à B4), \texttt{apc\_verdicts.json} (APC),
\texttt{coh\_verdicts.json} (C1, C2 ; le $\Delta$ de C1 contre la tendance n'est écrit que dans ce
fichier). Le $\Delta$ vs A0 est exploratoire : le critère N1 se juge contre la meilleure naïve. Les six
comparaisons contre A0 de \texttt{verdicts\_prive.json} ne sont pas corrigées pour la multiplicité ; avec
Bonferroni (IC à $1 - 0{,}05/6$), A2 (\enquote{bat A0}) et A1a (\enquote{échoue de peu}), ainsi que B4
(non candidat), passeraient à égalité\derive{} (approximation normale des IC à 95\,\%).}
\label{met:tab:ref14}
\end{table}

Toutes les victoires des variantes stock-flux contre la naïve tiennent à l'IC 99\,\% (Bonferroni, cinq
variantes ; \texttt{ci99\_bonferroni} de \texttt{verdicts\_prive.json})\mesure ; cet IC n'existe pas pour
l'APC, C1 et C2 ; à $h = 1$,
aucune variante n'est à plus d'un point de la persistance\mesure ; le banc des variantes stock-flux prend
18,7\mesure{}~s machine (\texttt{2026-11-14-axe2-methodes.md}). C1 bat A0 dans 17 cellules sur
22\mesure ; sans les deux cellules vues par le pilote, 7,52 contre 8,52\,\%\mesure
(\texttt{2026-11-14-axe2-cohortes-bayesiennes.md}). Par profession, C1 bat A0 pour les CD
($-0{,}92$ $[-1{,}45\,;\,-0{,}44]$\mesure), les IDE ($-0{,}96$ $[-1{,}48\,;\,-0{,}43]$\mesure) et les SF
($-1{,}99$ $[-3{,}83\,;\,-0{,}45]$\mesure), et fait jeu égal pour les MK (une origine) et les PH. Le gain de
C1 sur A0 décroît avec l'origine pour les CD et les PH, jusqu'à se renverser : à l'origine 2021, C1 fait
moins bien qu'A0 pour les CD ($+0{,}43$\mesure{} point) et les PH ($+1{,}24$\mesure{} point)
(\texttt{2026-11-14-axe2-cohortes-bayesiennes.md} l.~55--58). Sur le périmètre ancien (libéral + mixte), A0 mesure 8,21\,\%\mesure{} et A2 reste le seul gain
($-0{,}10$ $[-0{,}21\,;\,-0{,}002]$\mesure) ; B4 s'y renverse ($+0{,}53$ $[+0{,}15\,;\,+0{,}88]$\mesure)
(\texttt{verdicts\_ancien.json}).

\paragraph{Mondes simulés.} Sur 18 mondes à âges simples connus\suppose, l'erreur absolue sur l'effectif
des 60+ à $h = 5$ vaut 7,0\mesure{} (max. 18,9) pour le vieillissement uniforme, 5,3\mesure{} pour A1a et
2,1\mesure{} pour A1b ; sur le total, 3,9 / 3,0 / 3,8\,\%\mesure{} (\texttt{v2/methodes/out/monde.csv},
\texttt{2026-11-14-axe2-methodes.md}). A1b est donc retenu pour la structure par âge, sans gain sur le
total.

\paragraph{Prévision probabiliste (N6).} À l'origine 2021 (279 couples), le CRPS sur $\log N$ de C1 vaut,
sur la table du 14, 0,0654\mesure{} contre 0,0623\mesure{} pour la persistance probabiliste, $+0{,}003$
$[-0{,}001\,;\,+0{,}006]$\mesure{} (\texttt{coh\_verdicts.json}) ; sur la table du 15, 0,0656\mesure{}
contre 0,0625\mesure, $+0{,}003$ $[-0{,}001\,;\,+0{,}007]$\mesure{} (\texttt{c1i\_verdicts.json}) :
égalité dans les deux cas. \textbf{Ce n'est pas un résultat hors échantillon} : deux des trois cellules
jugées (CD et SF, origine 2021) ont été vues par le pilote avant la correction de la forme des variances
(\S\ref{met:bayes}) ; la couverture et le CRPS de C1 en N6 en héritent. Sa bande à 80\,\% couvre 73,5\,\%\mesure{} (cible
$[75\,;\,85]$), celle à 50\,\% 45,9\,\%\mesure{} (cible $[44\,;\,56]$) ; sur toutes les origines, hors
protocole, 77,6 et 50,0\,\%\mesure. L'APC seul couvre 50\,\%\mesure{} à 80\,\% nominal (PIT en U), parce
que ses entrées sont fixes (\texttt{coh\_verdicts.json}, \texttt{apc\_verdicts.json}). La calibration par
simulation de l'étape 2 (5 professions $\times$ 2 sexes $\times$ 200 jeux) échoue sur le niveau national
des entrées $L_T$ dans 10 cellules sur 10\mesure{} (\texttt{v2/methodes/out/coh\_sbc.json}). Aucune bande de C1 n'est calibrée : la bande nationale échoue à
la SBC, et la bande départementale échoue au critère de couverture de N6 (73,5\,\% $< 75$\,\%,
\texttt{calibre80} = \texttt{false}\mesure). La phrase de \texttt{decisions.md} selon laquelle \enquote{seule
la bande départementale de C1 est publiable} est donc inexacte.

\paragraph{Profil par horizon.} Le tableau~\ref{met:tab:horizon} donne la MAPE poolée par horizon du banc
du 15 (\texttt{c1\_gbm\_verdicts.json}, clé \texttt{mape\_par\_h}). C1 n'y figure pas : le budget d'erreur
ne l'a recalculé qu'à $h = 5$.

\begin{table}[ht]
\centering\small
\begin{tabular}{@{}lrrrrr@{}}
\toprule
$h$ (années) & 1 & 3 & 5 & 8 & 10 \\
\midrule
A2 & 2,73 & 5,59 & 8,60 & 14,33 & 17,59 \\
GBM global & 2,75 & 5,91 & 10,15 & 25,50 & 28,87 \\
MLP global & 3,70 & 9,77 & 20,52 & 23,33 & 51,15 \\
persistance & 3,12 & 6,85 & 9,94 & \textbf{13,56} & \textbf{16,06} \\
tendance & 3,70 & 6,39 & 9,32 & 14,42 & 17,78 \\
\bottomrule
\end{tabular}
\caption{MAPE poolée par horizon (\%), table privée du 2026-11-15, mesurée\mesure. Le nombre de cellules
décroît avec $h$ (les origines tardives n'ont pas de cible à 8 ou 10 ans).}
\label{met:tab:horizon}
\end{table}

\paragraph{GBM contre C1 sur le même millésime.} Le rapport \texttt{2026-11-15-c1-gbm-global.md} donne
$\Delta(\text{GBM} - \text{C1}) = +2{,}42$\mesure, IC 95\,\% $[+1{,}49\,;\,+3{,}70]$\mesure, mais en appariant
les erreurs du GBM (table du 15) à celles de C1 du 14 (\texttt{coh\_erreurs.parquet}) : deux millésimes de la
table SF. Sur le même millésime (tableau~\ref{met:tab:ref15}), la différence se dérive en point à partir des
MAPE par profession :
\begin{center}\small
\begin{tabular}{@{}lrrrrrr@{}}
\toprule
& Total & CD & IDE & MK & PH & SF \\
\midrule
$\Delta(\text{GBM} - \text{C1})$, même millésime & $+2{,}39$\derive & $+0{,}42$\derive & $+2{,}34$\derive & $+4{,}86$\derive & $+1{,}09$\derive & $+5{,}29$\derive \\
$\Delta(\text{A2} - \text{C1})$, même millésime & $+0{,}84$\derive & $+0{,}63$\derive & $+0{,}79$\derive & $+0{,}11$\derive & $+0{,}15$\derive & $+1{,}91$\derive \\
\bottomrule
\end{tabular}
\end{center}
\textbf{L'IC de ces différences n'est pas dérivable des sorties versionnées} : il faudrait les erreurs par
couple département $\times$ cellule de C1 sur la table du 15, que \texttt{methodes.budget\_erreur} écrit dans
\texttt{budget\_erreurs.parquet}, fichier ignoré par git et absent de \texttt{v2/methodes/out/} ; de même
\texttt{c1\_gbm\_erreurs.parquet} est absent. Le point $+2{,}39$ est à $0{,}03$ du $+2{,}42$ non apparié,
l'écart attendu d'après le seul changement de la table SF\derive. Le verdict (GBM rejeté contre C1) ne
dépend pas de ce choix ; la seule profession où il pourrait se discuter est CD ($+0{,}42$\derive, IC non
apparié $[-0{,}17\,;\,+1{,}01]$\mesure).

\subsection{Budget d'erreur}
\label{met:budget}
Code \texttt{v2/methodes/budget\_erreur.py} ; sorties \texttt{budget\_erreur.json} et
\texttt{budget\_erreur.log} ; rapport \texttt{2026-11-15-budget-erreur-c1-a2.md}. Même protocole que N1
(privé, $h = 5$, 2\,056 couples, mêmes tirages bootstrap) ; environ 6\mesure{} min machine.

\paragraph{Décomposition.} Avec des tranches de 5 ans et $h = 5$, la tranche $b$ à la cible vient exactement
de la tranche $b - 5$ à l'origine. L'erreur relative signée $e = (P - A)/A$ du département se coupe sans
reste en trois blocs selon l'âge à la cible, $e = e_Y + e_M + e_R$ ($Y$ : moins de 35 ans, surtout des
entrées ; $M$ : 35 à 59 ans, mobilité et sorties précoces ; $R$ : 60 ans et plus, retraites). Chaque $e_k$
se coupe en une part nationale (le total national du bloc par sexe $\times$ tranche) et une part locale (la
répartition entre départements). Les \emph{oracles} donnent au modèle la vraie valeur d'une composante à la
fois et mesurent le gain de MAPE. Les gains ne s'additionnent pas (la MAPE porte sur une valeur absolue) ; les
parts d'erreur quadratique $\mathrm E[e_k e]/\mathrm E[e^2]$, elles, somment à 1.

\paragraph{Plancher poissonien.} Si les événements individuels sont indépendants,
$\mathrm{Var}(N_T) \approx F$, avec $F = \sum |\text{solde de cohorte}|$ (borne basse des flux bruts), d'où
$\mathrm E|N_T - \mathrm E N_T| \approx \sqrt{2F/\pi}$ ; aucun prédicteur ponctuel ne descend sous ce plancher.

\begin{table}[ht]
\centering\small
\begin{tabular}{@{}lrr@{}}
\toprule
Points de MAPE (IC 95\,\%) & A2 & C1 \\
\midrule
MAPE $h = 5$\mesure & 8,60 $[7{,}80\,;\,9{,}47]$ & 7,76 $[7{,}12\,;\,8{,}40]$ \\
Plancher poissonien\derive & 3,11 $[2{,}89\,;\,3{,}36]$ & 3,11 $[2{,}89\,;\,3{,}36]$ \\
Part du plancher\derive & 36\,\% $[34\,;\,39]$ & 40\,\% $[37\,;\,43]$ \\
\addlinespace
Gain oracle $Y$ exact ($<35$)\mesure & 2,94 $[2{,}49\,;\,3{,}44]$ & 2,84 $[2{,}41\,;\,3{,}26]$ \\
\quad dont $Y$ local & 2,60 $[2{,}16\,;\,3{,}05]$ & 2,60 $[2{,}20\,;\,3{,}05]$ \\
\quad dont $Y$ national & 0,35 $[0{,}03\,;\,0{,}66]$ & 0,23 $[0{,}03\,;\,0{,}44]$ \\
Gain oracle $M$ exact (35--59)\mesure & 2,64 $[2{,}13\,;\,3{,}18]$ & 2,56 $[2{,}23\,;\,2{,}89]$ \\
\quad dont $M$ local (borne haute de la mobilité) & 2,00 $[1{,}58\,;\,2{,}46]$ & 2,36 $[2{,}00\,;\,2{,}71]$ \\
\quad dont $M$ national & 0,65 $[0{,}44\,;\,0{,}85]$ & 0,20 $[0{,}09\,;\,0{,}30]$ \\
Gain oracle $R$ exact (60+)\mesure & 0,23 $[0{,}08\,;\,0{,}40]$ & 0,22 $[0{,}07\,;\,0{,}37]$ \\
Gain oracle total national exact\mesure & 1,08 $[0{,}67\,;\,1{,}47]$ & 0,49 $[0{,}20\,;\,0{,}77]$ \\
\addlinespace
Écart local réduit de 25\,\% : $Y$ / $M$\mesure & 0,76 / 0,63 & 0,70 / 0,71 \\
Écart local réduit de 50\,\% : $Y$ / $M$\mesure & 1,37 / 1,18 & 1,29 / 1,30 \\
\addlinespace
Erreur quadratique totale (\%$^2$)\mesure & 153,0 & 117,1 \\
Part locale ($Y$ + $M$) de l'erreur quadratique\mesure & 77\,\% & 84\,\% \\
Part du bruit poissonien dans l'erreur quadratique\mesure & 12,6\,\% & 16,4\,\% \\
\bottomrule
\end{tabular}
\caption{Budget d'erreur au banc N1, table privée du 2026-11-15 (\texttt{budget\_erreur.json},
\texttt{budget\_erreur.log}). Parts locales : $0{,}413 + 0{,}357$ (A2), $0{,}446 + 0{,}395$ (C1)\mesure ;
le rapport arrondit celle de C1 à 85\,\%.}
\label{met:tab:budget}
\end{table}

\paragraph{Persistance de l'écart local.} Entre deux fenêtres disjointes (origines 2016 et 2021,
273\mesure{} couples département $\times$ profession, kinés exclus), la corrélation de l'écart local vaut,
pour A2 et C1 : $Y$ 0,17 et 0,02\mesure ; $M$ 0,24 et 0,22\mesure ; $R$ $-0{,}08$ et 0,16\mesure{} (erreur
type $\approx 0{,}06$\derive). Un correcteur fondé sur l'écart passé réduirait l'écart local d'au plus
$1 - \sqrt{1 - r^2} \le 3\,\%$\derive.

\paragraph{Selon la taille.} La MAPE de C1 vaut 14,04 / 7,60 / 5,26 / 5,49\,\%\mesure{} pour les
départements de 30--99, 100--299, 300--999 et 1\,000+ professionnels, pour un plancher de 6,29 / 3,42 /
1,92 / 1,07\derive. La pente de $\log$ APE sur $\log N$ vaut $-0{,}27$\mesure{} (A2 : $-0{,}29$\mesure)
contre $-0{,}54$\derive{} pour le plancher (environ $-1/2$ attendu) : les grands départements sont à environ
5 fois\derive{} leur plancher.

\paragraph{Par profession (C1).} SF 11,96 (plancher 4,86) ; CD 8,03 (3,14) ; MK 6,51 (1,93) ; IDE 4,93
(1,76) ; PH 4,91 (2,22)\mesure. Chez les SF, l'oracle national total gagne encore 1,07\mesure{} point ; chez
les CD, l'essentiel est $Y$ local (4,04\mesure).

\paragraph{Lecture.}
\begin{enumerate}
  \item L'hypothèse \enquote{le bruit imprévisible domine} est rejetée sous sa forme forte (plancher 40\,\%
    $<$ seuil de 50\,\%\suppose) et tenue sous sa forme faible : le reste est local et peu persistant
    ($r \le 0{,}24$\mesure). Ce rejet ne vaut que si les flux bruts égalent les soldes nets : le plancher
    croît comme $\sqrt F$, et le verdict cesse dès que les flux bruts dépassent environ
    $(0{,}5 \times 7{,}76/3{,}11)^2 \approx 1{,}6$\derive{} fois les soldes.
  \item Le niveau national explique la plus grande partie de l'écart C1 $-$ A2. Les gains d'oracle ne sont
    pas additifs, mais avec le total national exact pour les deux modèles, l'écart C1 $-$ A2 passe de
    $0{,}84$ à $(8{,}60 - 1{,}08) - (7{,}76 - 0{,}49) = 0{,}25$\derive{} point.
  \item Pour la mobilité inter-départementale, le critère (a) \enquote{gain réaliste $\ge 0{,}4$ point} est
    tenu (0,6 à 0,7\mesure) ; le critère (b) \enquote{persistance $r \ge 0{,}3$} n'est pas établi :
    $r$ = 0,22--0,24\mesure, IC 95\,\% $\approx [0{,}10\,;\,0{,}35]$\derive{} (transformation de Fisher,
    273 couples), qui contient 0,3 ; la persistance du signal n'est ni démontrée ni exclue. Conclusion de \texttt{decisions.md} : pas de modèle gravitaire, pas d'agents, pas
    de réseau profond. Les deux critères ont été posés après un premier passage, et git ne prouve
    pas leur antériorité (le code et les sorties entrent dans af0b816 à 01:50:25\mesure, le rapport et la
    décision dans 60d7a27 une seconde plus tard\mesure) : la conclusion est exploratoire.
\end{enumerate}

\subsection{Résultats négatifs et pistes écartées}
\label{met:negatifs}
\begin{description}
  \item[Intervalles conformes et MinT.] Test principal (origine 2021, 279 couples ; test d'abandon : CRPS inférieur à
    celui de C1). Le rapport dit ce test écrit avant le calcul, mais git ne le prouve pas : dérivation, code et
    résultats entrent dans le même commit 04eed25 (01:49:44\mesure). Les deux cellules CD et SF 2021 ont
    été vues par le pilote de C1 : le test n'est pas hors échantillon. Aucun des sept candidats (C1-conf,
    C1-conf-h5, C1-conf-hg, C1-MinT, C1-MinT-conf, C1-MinT-conf-hg, C1-MinT-conf-h5) ne passe : $\Delta$
    CRPS contre C1 de $+0{,}0012$ (C1-conf, $[-0{,}0014\,;\,+0{,}0040]$) à $+0{,}0098$ (C1-conf-h5,
    $[+0{,}0043\,;\,+0{,}0154]$)\mesure{} (\texttt{c1i\_verdicts.json}). Aucun ne bat la persistance
    probabiliste ; C1-conf-h5, C1-MinT et les trois C1-MinT-conf perdent contre elle\mesure. Le conforme
    abaisse la couverture à 80\,\% de 73,5\,\% à 62--67\,\%\mesure. Trois écarts contre C1 sont
    significatifs : MinT, $+0{,}0019$ $[+0{,}0010\,;\,+0{,}0030]$\mesure ; C1-conf-h5, $+0{,}0098$
    $[+0{,}0043\,;\,+0{,}0154]$\mesure ; C1-MinT-conf-h5, $+0{,}0092$ $[+0{,}0039\,;\,+0{,}0146]$\mesure. Sur toutes les origines (1\,581 couples,
    hors protocole), tous font jeu égal avec la persistance passée au même conforme\mesure. Cause
    mesurée : l'erreur de C1 n'est pas stationnaire. L'écart-type de $z$ croît de 0,80 à $h' = 1$ à 1,12
    à $h' = 5$\mesure ; la moyenne de $z$ à $h = 5$ passe de $-0{,}30$ (origine 2016) à $+0{,}59$ (2021)\mesure.
    Aucun recalage qui n'utilise que des cibles antérieures à l'origine ne voit ce changement. La
    couverture à 80\,\% de C1 en 2021 dépend de la taille : 91,4 / 74,2 / 54,8\,\%\mesure{} (petits,
    moyens, grands ; \texttt{c1i\_verdicts.json}). Décision (\texttt{decisions.md}) : la bande affichée ne change pas
    (c'est en fait $\pm k_p h$, \S\ref{met:servi}, et non celle des tirages de C1) ; le prochain effort porte sur le centre (choc commun national), pas sur la largeur
    (\texttt{decisions.md}).
  \item[GBM et MLP globaux.] Rejetés contre A2 (tableau~\ref{met:tab:ref15}) et contre C1 (même millésime :
    $+2{,}39$ et $+12{,}76$\derive{} point). Le GBM se dégrade avec l'horizon (28,87\,\%\mesure{} à
    $h = 10$, contre 17,59 pour A2). Biais national à $h = 5$ du GBM : SF $+14{,}24$\,\%\mesure, IDE
    $+8{,}52$\,\%\mesure. Taille d'échantillon : 32\,240 à 80\,620\mesure{} lignes, mais 4 à 9\mesure{}
    transitions annuelles et environ 100 départements $\times$ 5 professions de séries indépendantes\derive.
    Aucun réglage n'a été fait, et le référent inscrit était A2, pas C1 (\texttt{2026-11-15-c1-gbm-global.md},
    encadré de correction). L'antériorité du test d'abandon est prouvée par git : la dérivation entre dans
    70f8bc9 à 01:38:59\mesure, avant les résultats (d2819c6, 01:50:31\mesure). Ils restent au registre comme diagnostics.
  \item[Gravitaire et multi-agents.] Écartés par le budget d'erreur : la part locale de $M$ (borne haute de
    la mobilité) vaut 2,00 à 2,36\mesure{} point en oracle, mais l'écart local persiste peu
    ($r$ = 0,22--0,24\mesure{}, IC $\approx [0{,}10\,;\,0{,}35]$\derive, qui contient le seuil de 0,3) et aucun flux origine--destination ouvert n'existe. Les agents à
    règles calées étaient déjà tués a priori (non identifiables sans panel individuel).
  \item[Persistance meilleure qu'A2 à $h \ge 8$.] À $h = 8$ et $h = 10$, la persistance
    (13,56 et 16,06\,\%\mesure) fait mieux qu'A2 (14,33 et 17,59\,\%\mesure), et la tendance (14,42 et
    17,78) en est proche (tableau~\ref{met:tab:horizon}). Aucun IC n'a été calculé à ces horizons, et le
    critère N1 ne porte que sur $h = 5$ ; aucun rapport ne le commentait. L'affirmation \enquote{le
    stock-flux bat les références naïves} ne vaut donc qu'à cinq ans. C1 n'a pas de profil par horizon
    sur la table du 15.
  \item[Variantes stock-flux.] A3 dégrade la prévision ($+0{,}33$ $[+0{,}23\,;\,+0{,}42]$\mesure{} contre
    A0 ; biais national PH $-0{,}15 \to +2{,}28$\,\%\mesure, prédiction écrite falsifiée). A1a échoue de
    peu ($+0{,}09$ $[+0{,}01\,;\,+0{,}18]$\mesure) et coûte $+1{,}23$\mesure{}~s par profession à la
    commune. A123 n'apporte rien ($+0{,}06$ au privé, $+0{,}50$ $[+0{,}35\,;\,+0{,}64]$\mesure{} à
    l'ancien). B4 se renverse d'un périmètre à l'autre et produit une croissance composée
    (CD 2050/2026 : 1,61 contre 1,31\mesure{} sous A0, $+23\,\%$\derive).
  \item[C5, région de formation.] La part du stock 2026 prédit la part régionale des entrées nettes
    ($r = 0{,}996$\mesure) mieux que la part étudiante ($r$ = 0,64 CD, 0,94 PH\mesure).
  \item[C2.] Propager $\theta$ ne change pas la MAPE et dégrade légèrement la couverture (77,6 $\to$
    76,0\,\%\mesure) pour un coût double.
  \item[Échelle communale.] K : bande trop large (IDE 90 à 100\,\%\mesure) ; K2 : trop étroite au-dessus de
    200 professionnels (IDE 62,0, MK 39,6, CD 38,1\,\%\mesure) ; K3, K4, K5 : aucun grain
    infra-départemental calibré (K5 : IDE 65,9\,\% au bassin de vie et 59,1\,\% à l'EPCI\mesure).
    L'affichage reste au département (\texttt{2026-11-14-axe2-communes.md}).
\end{description}

\subsection{Limites}
\label{met:limites}
\begin{enumerate}
  \item \textbf{Cibles contaminées.} Les cibles 2023 et 2026 ont été vues par le banc K1 de la v1 ; le pilote
    de l'étape 2 a vu les cibles 2026 (origine 2021) des CD et des SF avant l'inscription de C1, puis la
    forme des variances a été corrigée. Sans ces deux cellules, C1 reste devant A0 (7,52 contre
    8,52\,\%\mesure), mais le seul test propre est le score des dépôts 2027 (temps du monde : RPPS 2027,
    vers juillet 2027\suppose).
  \item \textbf{IC trop étroits.} Le même tirage de départements sert pour toutes les cellules, ce qui
    conserve la corrélation entre origines d'un même département. Mais les origines (6 au plus) ne sont pas
    rééchantillonnées, le choc commun à tous les départements d'une même origine est ignoré, et la
    corrélation entre départements aussi\suppose. Remède : rééchantillonner aussi les origines, ou un
    bootstrap à deux grappes (départements et origines). La borne haute d'A2 contre A0 est à $-0{,}02$ (privé) et $-0{,}002$
    (ancien)\mesure. Les 2\,056 couples du CRPS ne sont pas indépendants.
  \item \textbf{Une seule origine jugeable en N6} (2021 ; SF, CD, PH ; 279 couples), dont deux cellules sur
    trois (CD et SF) vues par le pilote de C1. MK et IDE ne sont pas jugeables en probabiliste ; un seul choc
    commun est observé.
  \item \textbf{Biais non stationnaire de C1} : $z$ moyen de $-0{,}30$ à $+0{,}59$ entre les origines 2016 et
    2021\mesure ; gain sur A0 décroissant pour CD et PH, et renversé à l'origine 2021 (CD $+0{,}43$, PH
    $+1{,}24$\mesure{} point contre A0).
  \item \textbf{Millésimes mêlés.} La table privée SF a changé entre les bancs du 14 et du 15 ; toute
    comparaison entre les tableaux~\ref{met:tab:ref14} et~\ref{met:tab:ref15} mélange deux données. Relancer
    \texttt{methodes.banc} et \texttt{methodes.banc\_cohortes} sur la table actuelle coûte environ
    10\suppose{} min machine (\texttt{2026-11-15-budget-erreur-c1-a2.md}) ; ce n'est pas fait.
  \item \textbf{Périmètre privé reconstruit.} Les salariés ambigus sont répartis par une part FINESS 2026
    appliquée à toute la série ; seule la variante centrale est passée ; la série privée des SF est
    multipliée par 1,8\mesure{} entre 2012 et 2026. Les IDE sont mesurés au banc sur les libéraux RPPS, alors
    que la carte sert la table privée de l'axe 4, non mesurée ; série arrêtée en 2023.
  \item \textbf{Kinés : une seule origine} à $h = 5$ (IC de C1 contre A0 $[-1{,}39\,;\,+1{,}19]$\mesure ; IC
    du GBM contre A2 $[+1{,}13\,;\,+10{,}92]$\mesure).
  \item \textbf{Hypothèses d'âge supposées} : tranche $<25$ = 20--24 ans, 65+ de largeur 10, entrées
    uniformes ; à la commune, $<40$ = 25--39 ans\suppose.
  \item \textbf{Calibration par simulation partielle.} Pour l'APC à l'étape 1, un seul plan et 40 jeux avec
    $\theta$ ré-estimé ; à l'étape 2, plan réel seulement ; $L_T$ échoue (10/10\mesure), l'écart régional
    $\delta$ dans 3 cellules sur 10\mesure{} et l'écart départemental $\zeta$ dans 1 sur 10 (kinés,
    hommes)\mesure{} (\texttt{coh\_sbc.json}, seuil de Bonferroni).
  \item \textbf{Plafond $s = \min(1, e^{\eta})$} : biais de moyenne de $-0{,}55\,\%$\mesure{} au national
    (CD, origine 2021).
  \item \textbf{Budget d'erreur.} Le plancher est une borne basse (soldes nets, pas flux bruts) : avec des flux
    bruts deux fois plus grands, il monterait vers 4,4\suppose{} points, 57\,\%\derive{} de la MAPE de C1, et
    le test d'abandon basculerait. Le découpage par âge est un proxy (les entrées tardives tombent dans $M$).
    La persistance repose sur une seule paire d'origines, sans les kinés. L'oracle national garde la
    répartition du modèle. Le champ \texttt{"cellules": 11} de \texttt{budget\_erreur.json} est faux (division
    par le nombre de modèles, \texttt{budget\_erreur.py}, fonction \texttt{main}) ; le banc a bien 22 cellules.
  \item \textbf{GBM.} Aucun réglage, qui pourrait réduire l'écart ; un essai de fumée a vu une cible de test
    (2021 $\to$ 2026) ; modèle direct à $h = 5$ et GBM sur résidus d'A2 non testés.
  \item \textbf{Intervalles.} Références de méthode non lues en texte intégral ; C1-conf-hg post hoc ; MinT
    avec région et France par relabellisation, ce qui réduit le rétrécissement régional à un seul groupe.
  \item \textbf{Jumeau.} C1 est injecté par un cache de l'historique, un détail interne du jumeau ; sa
    rétention est tenue constante sur l'horizon ($-1{,}1$ à $+1{,}2\,\%$\mesure{} sur le total national à
    5 ans, \texttt{decisions.md}) ; le levier \texttt{decalage\_retraite} hérite d'un biais sur la tranche 65+.
    Deux mesures de recalcul coexistent : 14 à 40\mesure{}~ms (A0 et variantes,
    \texttt{v2/papers/axe2-methodes/main.tex}) et 280 à 330\mesure{}~ms (C1 servi, tous leviers,
    \texttt{decisions.md}).
  \item \textbf{Au-delà de $h = 10$}, rien n'est rétro-testable ; les projections à 2050 (A0 contre B4 contre
    la DREES~\cite{b-drees2021dd76}) ne sont pas départagées (protocole N2 non passé).
\end{enumerate}
La revue systématique de~\cite{b-lee2024hrh} note que la validation contre des données passées reste rare
dans les modèles de projection d'effectifs ; le banc N1 en est une, avec les limites ci-dessus. La forme
bayésienne par composantes de C1 est voisine de~\cite{b-alexander2021bayesian}, sans la contrainte de somme
au total national publié.
\endgroup
